\chapter{Bidirectional interaction: conjugation of advanced and retarded propagators and boundary conditions of the absorber}
\label{chap:mt-accion-interaccion-absorbedor}

A route becomes dynamical when it is endowed with an action functional and boundary conditions. The action assigns a weight to the complete history; its variation selects transports compatible with the endpoints; interaction couples degrees of freedom that share a boundary; the bidirectional reading preserves the forward and reverse temporal orientations of the same antecedent. Within this framework, emission and absorption form the endpoints of a common boundary-value problem.

The absorber occupies a precise position in the architecture. HMT supplies the
oriented band, outward-and-return memory, and the discrete involutions.
Dimensional Mechanics adds the evolution operator, currents, and the
action transducer. A physical absorber realization links that
structure to advanced and retarded propagators and causal boundary
conditions. The algebraic identities of the first segment are
proved in their domains; the final link retains the status of a
physical interface until the complete intertwiner is constructed.

\section{Action functional on enriched paths}

Let $\gamma=((c_0,\mu_0),(d_1,\ldots,d_L))$ be a finite route, with directions
$d_t$ and sheets $\mu_t\in\{\Sigma,\Pi\}$. Define
\[
 \operatorname{turn}(\gamma)
 =\#\{2\leq t\leq L:d_t\neq d_{t-1}\},
\]
\[
 \operatorname{switch}(\gamma)
 =\#\{2\leq t\leq L:\mu_t\neq\mu_{t-1}\}.
\]

\begin{definicionnarrativa}[title={Discrete Action of Minimal Interaction}]
\label{def:mt-accion-minima}
For $\lambda\geq0$, the local route action is
\begin{equation}
 \mathcal A_\lambda(\gamma)
 =\operatorname{turn}(\gamma)
 +\lambda\operatorname{switch}(\gamma)
 =\sum_{t=2}^{L}
 \bigl[1-\delta_{d_t,d_{t-1}}
 +\lambda(1-\delta_{\mu_t,\mu_{t-1}})\bigr].
 \label{eq:mt-accion-minima}
\end{equation}
The first contribution measures directional variation and the second measures transport between leaves.
\end{definicionnarrativa}

The general action also incorporates boundary, carry, and memory:
\begin{equation}
 \mathcal S[\gamma]
 =\sum_{e\in\gamma}\mathcal L_e
 +\mathcal B_{\rm in}(\partial_-\gamma)
 +\mathcal B_{\rm out}(\partial_+\gamma)
 +\mathcal M(\gamma).
 \label{eq:mt-accion-ruta-general}
\end{equation}
Here $\mathcal L_e$ is the local cost, $\mathcal B_{\rm in/out}$ fixes the
admissible endpoint data, and $\mathcal M$ preserves the information that
the visible publication contracts. The variation is performed within a class
of routes with the same boundary types.

\begin{lecturafisica}[title={Principle of Minimal Interaction Compatible with Memory}]
\label{prin:mt-minima-interaccion}
Among the prolongations compatible with the boundary data, the
selected physical realization extremizes the action functional within the
sector that preserves the memory required by the composition.
\end{lecturafisica}

The principle fixes the class of the variational problem. A concrete dynamic
adds the Hamiltonian or Lagrangian of the sector and proves that its extremum
produces the equations of motion. Memory functions as a dynamic constraint:
two routes with the same final shadow may belong to distinct
variational classes.

\section{Interaction as typed coupling}

Let $\mathfrak p_1$ and $\mathfrak p_2$ two sections realized over routes
$\gamma_1$ and $\gamma_2$. An interaction at to boundary $B$ consists of
to morphism
\begin{equation}
 \mathcal I_B:
 E^{(1)}_B\otimes E^{(2)}_B
 \longrightarrow E^{(12)}_B
 \label{eq:mt-interaccion-frontera}
\end{equation}
and in an action term $\mathcal S_{\rm int}$ compatible with the
representations of the sector. The morphism preserves the shared types and
transforms the internal degrees according to the coupling rules. The
complete composition takes the form
\[
 (\gamma_2,s_2)\star_{(B,\mathcal I_B)}(\gamma_1,s_1),
\]
where the record of $B$ is counted once, and the memory of both sides is
transported to the composite state.

The covariance of the interaction requires, for each symmetry transformation
$g$,
\begin{equation}
 \rho_{12}(g)\,\mathcal I_B
 =\mathcal I_B\bigl(\rho_1(g)\otimes\rho_2(g)\bigr).
 \label{eq:mt-covariancia-interaccion}
\end{equation}
This identity decides whether the coupling belongs to the declared
physical representation. In the electromagnetic and color sectors, the same
structure is expressed by means of a connection and its parallel transport.

\section{Temporal conjugation of the propagator}

Let $\mathcal H_{\mathbb R}$ the realification of the quantum space, endowed
with to complex structure $J_E$ with $J_E^2=-I$. Let $H=H^*$ to self-adjoint operator
with common invariant domain and is $C$ an orthogonal involution
that satisfies
\begin{equation}
 HJ_E=J_EH,\qquad
 CJ_EC^{-1}=-J_E,
 \qquad CHC^{-1}=H.
 \label{eq:mt-hipotesis-reversion-propagador}
\end{equation}

\begin{resultado}[title={Theorem: Bidirectional reversal of the evolution group}]
\label{thm:mt-reversion-propagador}
The operator
\begin{equation}
 U(t)=\exp\!\left(-\frac{J_EHt}{\hbar_{\rm int}}\right)
 \label{eq:mt-grupo-evolucion}
\end{equation}
forms a uniparametric unitary group and satisfies
\begin{equation}
 \boxed{CU(t)C^{-1}=U(-t)}.
 \label{eq:mt-reversion-unitaria}
\end{equation}
\end{resultado}

\paragraph{Proof of the bidirectional reversal of the unitary group \(U(t)\): \(CU(t)C^{-1}=U(-t)\)}
The generator $A=-J_EH/\hbar_{\rm int}$ is skew-adjoint. Stone's theorem
produces the unitary group $U(t)=e^{tA}$. The hypotheses of
\eqref{eq:mt-hipotesis-reversion-propagador} give $CAC^{-1}=-A$; functional
calculus yields $Ce^{tA}C^{-1}=e^{-tA}$.


The involution $C_M=-\operatorname{rev}$ of the route words supplies the
genealogical datum of orientation. When the physical functor carries $C_M$ to the
involution $C$ of the theorem, the particle and antiparticle sections are
represented by conjugate temporal orientations of the same unitary
group. The identity is proved from explicit starting structure:
self-adjointness, domain, complex structure, and involution.

\section{Advanced and Retarded Propagators}

Let $L$ be the hyperbolic operator of a field realization and let
$G_{\rm ret}$ and $G_{\rm adv}$ be its fundamental inverses with future and past
causal supports, respectively:
\begin{equation}
 LG_{\rm ret}=LG_{\rm adv}=I,
 \qquad
 \operatorname{supp}(G_{\rm ret}f)\subset J^+(\operatorname{supp}f),
 \qquad
 \operatorname{supp}(G_{\rm adv}f)\subset J^-(\operatorname{supp}f).
 \label{eq:mt-green-avanzado-retardado}
\end{equation}
The causal difference
$E=G_{\rm ret}-G_{\rm adv}$ transports data between solutions, and the symmetric
mean
\[
 G_{\rm sym}=\frac12(G_{\rm ret}+G_{\rm adv})
\]
realizes a time-symmetric kernel.

An HMT--MD interface to this theory is a family of maps
\begin{equation}
 \mathfrak I_{\rm prop}:
 (\gamma,\gamma^\vee,\mathcal B_\gamma,U,C_M)
 \longrightarrow
 (G_{\rm ret},G_{\rm adv},E,\mathcal H_{\rm field})
 \label{eq:mt-interfaz-green}
\end{equation}
which intertwines the route involution with the exchange
$G_{\rm ret}\leftrightarrow G_{\rm adv}$, preserves the inner product or the
relevant symplectic form, and carries the actions on both sides to the same
functional form. This map is the mathematical condition that turns the
internal bidirectional orientation into advanced--retarded propagation of
fields.

\section{HMT--MD realization of the absorber through route inversion and particle--antiparticle conjugation}

\begin{definicionnarrativa}[title={Bidirectional emission--absorption system}]
\label{def:mt-absorbedor}
A physical realization of an absorber is a boundary-value problem formed by
an emitting current $J_{\rm em}$, an absorbing response $J_{\rm abs}$,
the propagators $G_{\rm ret}$ and $G_{\rm adv}$, and boundary conditions that
jointly close the two orientations of transport. Its temporally symmetric
effective action has the form
\begin{equation}
 \mathcal S_{\rm abs}[J]
 =\mathcal S_{\rm matter}[J]
 +\frac12\langle J,G_{\rm sym}J\rangle,
 \qquad J=J_{\rm em}+J_{\rm abs},
 \label{eq:mt-accion-absorbedor}
\end{equation}
and its boundary conditions determine the observed radiative part.
\end{definicionnarrativa}

The definition expresses a theory of action at a distance or boundary response in operator language. Within HMT--MD, the band supplies the pair $(\gamma,\gamma^\vee)$; the bidirectional-reversal result supplies the exact reversal of the unitary group; the functor $\mathfrak I_{\rm prop}$ must supply causality, an inner product, and boundary conditions. The absorber architecture prolongs the bidirectional structure already constructed. Its realization remains a physical interface, and the construction of the complete causal intertwiner constitutes an explicitly delimited mathematical problem.

The physical closure requires four verifiable identities:
\begin{enumerate}
 \item commutativity between route conjugation and interchange of the two
 propagators;
 \item equality of the actions transported by
 $\mathfrak I_{\rm prop}$;
 \item preservation of the symplectic form or the inner product;
 \item compatibility of the emission and absorption conditions under
 route refinement.
\end{enumerate}
TO failure in any of these identities refutes the particular realization
and leaves intact the operatorial reversion already proved in
\eqref{eq:mt-reversion-unitaria}.

\section{The separate actions of \texorpdfstring{$C$}{C},
\texorpdfstring{$P$}{P} and \texorpdfstring{$T$}{T}}

On a complete route we define three typed involutions:
\begin{enumerate}
 \item $C$ exchanges the sheets $\Sigma\leftrightarrow\Pi$ and conjugates the
 oriented characters;
 \item $P$ reflects the spatial chart, for example
 $(i,j)\mapsto(i,-j)\pmod9$, together with
 $E\leftrightarrow W$;
 \item $T$ reverses the order of the record, replaces each direction with its
 opposite, and permutes the endpoints of the route.
\end{enumerate}
Each operation acts on to distinct component of the state. Their composition
$CPT$ acts on sheet, space, and temporal orientation at once.

\begin{resultado}[title={Theorem: Individual and joint discrete invariance}]
\label{thm:mt-cpt-rutas}
For every route $\gamma$ and every $\lambda\geq0$,
\begin{equation}
 \mathcal A_\lambda(\gamma)
 =\mathcal A_\lambda(C\gamma)
 =\mathcal A_\lambda(P\gamma)
 =\mathcal A_\lambda(T\gamma)
 =\mathcal A_\lambda(CPT\gamma).
 \label{eq:mt-invariancia-cpt-rutas}
\end{equation}
\end{resultado}

\paragraph{Compatibility of the Actions \(C\), \(P\), and \(T\) with the Unitary Reversal \(U(t)\mapsto U(-t)\)}
$P$ permutes the alphabet of directions and preserves the pattern of equalities between consecutive steps. $C$ globally permutes the sheets and preserves the pattern of sheet changes. $T$ reverses the order of both sequences and preserves the number of adjacent transitions. Each involution separately preserves the two summands of \eqref{eq:mt-accion-minima}; the composition inherits the invariance.


The theorem belongs to the discrete ledger and its local action. The CKM matrix realized in the flavor basis has a nonzero Jarlskog invariant and therefore CP violation in that representation. Both results are compatible: an oriented reader can distinguish states related by $CP$ while the local functional $\mathcal A_\lambda$ preserves each involution. The realization of the relativistic field-theoretic CPT theorem adds locality, Lorentz covariance, spectral structure, and a representation of fields. In that enlarged domain, joint CPT constitutes the symmetry to be preserved. In the internal domain, the CKM closure and the temporal expansion remain compatible with the CPT invariance proved for the discrete action.

Three strata are thereby kept separate. The invariance of
\eqref{eq:mt-invariancia-cpt-rutas} is an exact internal theorem. The identity
\eqref{eq:mt-reversion-unitaria} has been proved from explicit starting
structure. The physical absorber theory and the CPT theorem for fields require
the intertwiners proper to their codomains. This separation keeps the
inverse temporal orientation within the definition of a particle and accords
each assertion its appropriate force.
